\section{Attrazione utenti}
\label{sec:attrazione_utenti}

Oltre a descrizioni accattivanti nei meta tag, sono state adottate ulteriori strategie per attrarre e coinvolgere gli utenti o rispondere a specifiche esigenze.
In generale per rispondere a più tipologie di utenti abbiamo adottato schemi organizzativi ambigui, si rimanda alla sezione \ref{sec:architettura_informazione} per maggiori dettagli.
\paragraph{Pagina Instagram}
È stata creata una pagina Instagram dedicata al progetto per promuovere il rifugio e gli animali in cerca di adozione. La pagina condivide post di benvenuto per i nuovi ospiti e informazioni utili sulla struttura.

\paragraph{Sistema dei Preferiti}
La funzione dei preferiti offre agli utenti un modo rapido per salvare gli animali di interesse e riprendere la navigazione successivamente. Nella metafora della pesca, ciò corrisponde alla "boa di segnalazione" che aiuta a ritrovare i "pesci" catturati in precedenza.

Tuttavia, la funzione va oltre un semplice segnalibro: un utente non molto sicuro di voler adottare potrebbe inserire nei preferiti l'animale, con intenzioni non certe di ritornare al sito. Una call to action però, propone il login per sincronizzare i preferiti su tutti i dispositivi. In questo modo si crea un impegno psicologico che incentiva il ritorno sulla piattaforma.

\paragraph{Filtri di ricerca}
Per rispondere a più tipologie di utente, abbiamo implementato filtri di ricerca ambigui in alcune sezioni del sito. Ad esempio, nella pagina degli animali adottati, il filtro "Tutte" include sia gli animali assegnati all'amministratore che quelli non assegnati, permettendo agli utenti di esplorare una gamma più ampia di opzioni senza limitarsi a categorie specifiche.\\
Inotre sono utilizzati anche filtri esatti (es. ricerca per gatti o cani) o ricerca per nome per soddisfare utenti con esigenze più precise.
Questa strategia corrisponde all'idea di "pesca con l'amo" (o "tiro perfetto") nella metafora della pesca.

\paragraph{Home e struttura a ipertesti}
La struttura a ipertesti (descritta nella sezione \ref{sec:strutture_organizzative}) consente di guidare l'utente seguendo percorsi specifici, che lo aiutano ad avvicinarsi sempre di più all'informazione cercata.
Si può considerare la home page come una "tappola per aragoste" nella metafora della pesca: essa attira l'utente con informazioni generali e lo indirizza verso sezioni più specifiche del sito, facilitando la scoperta di contenuti rilevanti.
\paragraph{Sezione eventi}
Aprendo la scheda dettaglio di un evento, l'utente può visulizzare a lato altri eventi che si tengono nella stessa zona. Questo aiuta l'utene a trovare l'evento che più gli interessa in base alla sua posizione geografica, aumentando così la probabilità di partecipazione.
Nella metafora della pesca, ciò è descritto come "trappola per aragoste".